% GENERATED FILE. Edit canonical lessons, then run npm run generate:books.
% canonical-source-hash: fnv1a64:020da524b4dbca1e
% canonical-lessons: GE-C250-ausserdem, GE-C250-deshalb, GE-C250-trotzdem, GE-C250-jedenfalls, GE-C250-eigentlich

\chapter{Besides, Therefore, Nevertheless, In any case, Actually}
\label{ch:ge-besides-therefore-nevertheless-in-any-case-actually}
\hlchaptermodality{250}

\begin{chapteropening}
\textbf{By the end of this chapter:} \emph{I can say besides, therefore, nevertheless, in any case and actually.}

Complete the last lesson of chapter 250: I can say besides, therefore, nevertheless, in any case and actually.
\end{chapteropening}

\section[außerdem]{außerdem — besides, moreover}

\label{lesson:GE-C250-ausserdem}



\begin{quote}
Before the new one: say the German for to produce, then the German for to demand.
\end{quote}

\subsection*{You'll want to know: außerdem}
\textbf{außerdem} — \textquotedblleft{}besides, moreover\textquotedblright{}.

\begin{cousinweb}
\textbf{außerdem} is the word for \textquotedblleft{}besides, outside\textquotedblright{}, built on the word for \textquotedblleft{}out\textquotedblright{}, + \emph{dem}, a form of \emph{der}: \textquotedblleft{}besides that\textquotedblright{}.
\end{cousinweb}

\begin{grammarlens}[title={What you've built}]
One more word for this chapter.
\end{grammarlens}

\subsection*{Your turn}
\begin{itemize}
\raggedright
  \item \emph{Say it:} \emph{außerdem}
  \item \emph{Say it:} \emph{außerdem}, once more
  \item \emph{Say it:} say \emph{außerdem}
  \item \emph{Recall:} say the German for to produce, then the German for to demand, then say \emph{außerdem} again
\end{itemize}

\begin{tcolorbox}[breakable,colback=teal!4,colframe=teal!35!black,title={Before you move on}]
What does außerdem mean? (\textquotedblleft{}Besides, moreover\textquotedblright{}.) Say it once more.
\end{tcolorbox}
\section[deshalb]{deshalb — therefore}

\label{lesson:GE-C250-deshalb}



\begin{quote}
Before the new one: say the German for to demand, then the German for besides.
\end{quote}

\subsection*{You'll want to know: deshalb}
\textbf{deshalb} — \textquotedblleft{}therefore\textquotedblright{}.

\begin{cousinweb}
\textbf{deshalb} is \emph{des}, \textquotedblleft{}of that\textquotedblright{}, a form of \emph{der}, + \emph{halb}, an old word for \textquotedblleft{}side, behalf\textquotedblright{} (as in English \emph{behalf}): on that account.
\end{cousinweb}

\begin{grammarlens}[title={What you've built}]
One more word for this chapter.
\end{grammarlens}

\subsection*{Your turn}
\begin{itemize}
\raggedright
  \item \emph{Say it:} \emph{deshalb}
  \item \emph{Say it:} \emph{deshalb}, once more
  \item \emph{Say it:} say \emph{deshalb}
  \item \emph{Recall:} say the German for to demand, then the German for besides, then say \emph{deshalb} again
\end{itemize}

\begin{tcolorbox}[breakable,colback=teal!4,colframe=teal!35!black,title={Before you move on}]
What does deshalb mean? (\textquotedblleft{}Therefore\textquotedblright{}.) Say it once more.
\end{tcolorbox}
\section[trotzdem]{trotzdem — nevertheless}

\label{lesson:GE-C250-trotzdem}



\begin{quote}
Before the new one: say the German for besides, then the German for therefore.
\end{quote}

\subsection*{You'll want to know: trotzdem}
\textbf{trotzdem} — \textquotedblleft{}nevertheless\textquotedblright{}.

\begin{cousinweb}
\textbf{trotzdem} is the word for \textquotedblleft{}defiance\textquotedblright{}, Middle High German \emph{traz} (EtymDB 2.1), + \emph{dem}, a form of \emph{der}: in defiance of that.
\end{cousinweb}

\begin{grammarlens}[title={What you've built}]
One more word for this chapter.
\end{grammarlens}

\subsection*{Your turn}
\begin{itemize}
\raggedright
  \item \emph{Say it:} \emph{trotzdem}
  \item \emph{Say it:} \emph{trotzdem}, once more
  \item \emph{Say it:} say \emph{trotzdem}
  \item \emph{Recall:} say the German for besides, then the German for therefore, then say \emph{trotzdem} again
\end{itemize}

\begin{tcolorbox}[breakable,colback=teal!4,colframe=teal!35!black,title={Before you move on}]
What does trotzdem mean? (\textquotedblleft{}Nevertheless\textquotedblright{}.) Say it once more.
\end{tcolorbox}
\section[jedenfalls]{jedenfalls — in any case}

\label{lesson:GE-C250-jedenfalls}



\begin{quote}
Before the new one: say the German for therefore, then the German for nevertheless.
\end{quote}

\subsection*{You'll want to know: jedenfalls}
\textbf{jedenfalls} — \textquotedblleft{}in any case\textquotedblright{}.

\begin{cousinweb}
\textbf{jedenfalls} is the word for \textquotedblleft{}every\textquotedblright{} + the word for \textquotedblleft{}case\textquotedblright{}, from \emph{fallen}, + \emph{-s}: in every case.
\end{cousinweb}

\begin{grammarlens}[title={What you've built}]
One more word for this chapter.
\end{grammarlens}

\subsection*{Your turn}
\begin{itemize}
\raggedright
  \item \emph{Say it:} \emph{jedenfalls}
  \item \emph{Say it:} \emph{jedenfalls}, once more
  \item \emph{Say it:} say \emph{jedenfalls}
  \item \emph{Recall:} say the German for therefore, then the German for nevertheless, then say \emph{jedenfalls} again
\end{itemize}

\begin{tcolorbox}[breakable,colback=teal!4,colframe=teal!35!black,title={Before you move on}]
What does jedenfalls mean? (\textquotedblleft{}In any case\textquotedblright{}.) Say it once more.
\end{tcolorbox}
\section[eigentlich]{eigentlich — actually}

\label{lesson:GE-C250-eigentlich}



\begin{quote}
Before the new one: say the German for nevertheless, then the German for in any case.
\end{quote}

\subsection*{You'll want to know: eigentlich}
\textbf{eigentlich} — \textquotedblleft{}actually\textquotedblright{}.

\begin{cousinweb}
\textbf{eigentlich} $\leftarrow$ Middle High German \emph{eigenlīche} (EtymDB 2.1): the word for \textquotedblleft{}own\textquotedblright{}, Old High German \emph{eigan}, English \emph{own}, + \emph{-lich}: in its own nature, actually.
\end{cousinweb}

\begin{grammarlens}[title={What you've built}]
One more word for this chapter.
\end{grammarlens}

\subsection*{Your turn}
\begin{itemize}
\raggedright
  \item \emph{Say it:} \emph{eigentlich}
  \item \emph{Say it:} \emph{eigentlich}, once more
  \item \emph{Say it:} say \emph{eigentlich}
  \item \emph{Recall:} say the German for nevertheless, then the German for in any case, then say \emph{eigentlich} again
\end{itemize}

\begin{tcolorbox}[breakable,colback=teal!4,colframe=teal!35!black,title={Before you move on}]
What does eigentlich mean? (\textquotedblleft{}Actually\textquotedblright{}.) Say it once more.
\end{tcolorbox}
